\documentclass[12pt]{article}

\usepackage[T1]{fontenc}
\usepackage[utf8]{inputenc}
\usepackage[brazil]{babel}
\usepackage{graphicx}
\usepackage{url}
\usepackage{sbc-template}

\sloppy

\title{Análise Empírica e Otimização para o Problema do Caixeiro Viajante}

\author{Nome do(s) autor(es)}

\address{Departamento de Ciência da Computação -- Universidade Federal de Minas Gerais
  \email{email@ufmg.br}
}

\begin{document}

\maketitle

\begin{abstract}
% To be written after the experiments are completed.
\end{abstract}

\begin{resumo}
% A ser escrito após a conclusão dos experimentos.
\end{resumo}

\section{Introdução}
% Contexto, problema, objetivos e organização do artigo.

\section{Implementação, Profiling e Auditoria}
% Estruturas de dados, decisões algorítmicas, refatorações e gargalos.

\subsection{Branch-and-Bound e estratégias de warm-start}

\subsection{Algoritmos aproximativos}

\subsection{Instrumentação e profiling}

\section{Experimentos e Resultados}
% Não inserir resultados antes da execução e análise dos experimentos.

\subsection{Modelagem de complexidade empírica}

\subsection{Impacto do warm-start na árvore de busca}

\subsection{Análise de sensibilidade das instâncias}

\section{Conclusões}
% Síntese das descobertas, limitações e considerações finais.

\end{document}
